% !TeX root = ../../Book2.tex
% 纲要标题：### C.4 向量丛、自同态代数与 Morita 等价
% 纲要位置：工作记录/计划纲要.md，第 1828 行。

% 纲要要点（供目录核对）：
% * 有限投射模的局部坐标、自同态代数的粘合与 Morita 等价
% #### C.4.1 向量丛与自同态代数
% #### C.4.2 模范畴与几何信息
% ---

\section{向量丛、自同态代数与 Morita 等价}
\label{b2:sec:C04}

有限生成投射模在每个素理想处给出一个有限维纤维；要描述这些纤维如何组成同一个对象，还需主开覆盖上的自由模，以及交叠处满足相容条件的换基矩阵。这个局部坐标描述把有限生成投射模与向量丛联系起来，并使其自同态代数成为由共轭变换粘合的局部矩阵代数。Morita 等价则说明，从模范畴观察这些自同态代数时，哪些来自交换基环的信息得到保留。

\subsection{向量丛与自同态代数}
\label{b2:subsec:C04:01}

设 \(R\) 为非零交换含幺环、\(P\) 为有限生成投射 \(R\) 模。对 \(f\in R\)，主开集 \(D(f)=\{\mathfrak p\in\operatorname{Spec}R\mid f\notin\mathfrak p\}\) 上的模由 \(P_f=R[1/f]\otimes_RP\) 描述。\cref{b2:prop:projective-basic-open-basis}给出有限主开覆盖 \(D(f_i)\)，使每个 \(P_{f_i}\) 都自由。若局部秩均为 \(n\)，在这组覆盖中分别选基；交叠区域上的坐标由换基矩阵
\[
g_{ij}\in\operatorname{GL}_n(R_{f_if_j})
\]
变换，并在三重交叠上满足 \(g_{ij}g_{jk}=g_{ik}\)。这里约定 \([p]_i=g_{ij}[p]_j\)，两边均取到 \(R_{f_if_j}\) 上。这些主开集上的自由模及其相容换基矩阵，给出\term[xiangliang-cong]{向量丛}[vector bundle]的局部描述；本节只使用从已经给定的 \(P\) 导出的这一模型。不同连通部分可以有不同的局部秩。

\ProseParagraphBreak
对 \(E=\operatorname{End}_R(P)\)，同一组局部基给出 \(E_{f_i}\cong M_n(R_{f_i})\)。若 \(T\in E\) 的局部矩阵为 \(T_i\)，由坐标变换公式可得
\begin{equation}\label{b2:eq:endomorphism-bundle-transition}
T_i=g_{ij}T_jg_{ij}^{-1}.
\end{equation}
所以自同态代数使用共轭来粘合矩阵代数。这些公式保留加法、乘法与单位矩阵；因此局部矩阵的描述与全局代数结构相容。若 \(P\) 忠实，它给出\cref{b2:cor:azumaya-endomorphism}中的分裂 Azumaya 代数。

\ProseParagraphBreak
局部秩为一时，\(g_{ij}\) 是标量单位，局部自同态 \(T_j\) 也是标量，交换性使 \(g_{ij}T_jg_{ij}^{-1}=T_j\)。因此模本身的换基可能仍然非平凡，自同态代数的共轭换基却已经消失。是否存在全局模基，与自同态代数是否为矩阵环，是两个不同的问题。

\ProseParagraphBreak
\cref{b2:ex:krull-schmidt-nonunique}中的 \(R=\mathbb Z[s]\)、\(s^2=-5\)、\(I=(2,1+s)\) 给出同一局部秩下的例子。该处已证 \(I\) 非主且 \(I^2=2R\)。对 \(t\in I\)，映射 \(t\mapsto2t\) 及 \(t\mapsto(s-1)t/2\) 都取值于 \(R\)，后者使用 \(s-1\in I\) 和 \(I^2=2R\)。并有
\[
t=2(2t)+(1+s)\frac{(s-1)t}{2}.
\]
由\cref{b2:prop:finite-dual-basis}，\(I\) 是有限生成投射模。在 \(D(2)\) 上，\(I_2=R_2\)；在 \(D(3)\) 上，由 \(2=(1+s)(1-s)/3\)，有 \(I_3=(1+s)R_3\)。这两个主开集覆盖 \(\operatorname{Spec}R\)，所以局部秩处处为一；若有全局基，\(I\) 就应由一个元素生成，与非主性矛盾。

\ProseParagraphBreak
尽管 \(I\) 不自由，却有自然的代数同构 \(\operatorname{End}_R(I)\cong R\)。确实，取 \(0\ne t_0\in I\)，任一自同态 \(T\) 满足 \(t_0T(t)=tT(t_0)\)，所以它在分式域中是乘以某个 \(c\) 的映射。由 \(cI\subseteq I\)，有 \(cI^2\subseteq I^2=2R\)，消去 \(2\) 得 \(cR\subseteq R\)，故 \(c\in R\)。反之每个 \(c\in R\) 都给出这样的自同态，乘法和复合相容。这具体展示了秩一换基的共轭作用为何消失。

\ProseParagraphBreak
例如 \(R=k\times k\)，取 \(P=k\times k^2\)，由两个坐标分别作标量作用。这个模是 \(R^2\) 的直和因子且忠实，两个点上的秩分别为一、二，故不可能是某个固定秩的自由 \(R\) 模。其自同态代数为 \(k\times M_2(k)\)，是分裂 Azumaya \(R\) 代数，却不同构于任何 \(M_n(R)\)：在两个坐标上比较维数将同时要求 \(n^2=1\) 与 \(n^2=4\)。这正是环上的“分裂”比“一个固定大小的矩阵环”更自然的原因。

\subsection{模范畴与几何信息}
\label{b2:subsec:C04:02}

\begin{proposition}\label{b2:prop:split-azumaya-morita}
设 \(R\) 为非零交换含幺环、\(P\) 为忠实有限生成投射 \(R\) 模，令 \(E=\operatorname{End}_R(P)\)、\(P^*=\operatorname{Hom}_R(P,R)\)。赋予 \(P\) 双模结构 \({}_EP_R\)，赋予 \(P^*\) 双模结构 \({}_RP^*_E\)，其中
\[
ep=e(p),\qquad pr=rp,\qquad
(r\lambda)(p)=r\lambda(p),\qquad (\lambda e)(p)=\lambda(e(p)).
\]
则左模范畴之间的两张量函子
\[
\begin{aligned}
F:R\text{-}\mathbf{Mod}&\longrightarrow E\text{-}\mathbf{Mod},
&N&\longmapsto P\otimes_RN,\\
G:E\text{-}\mathbf{Mod}&\longrightarrow R\text{-}\mathbf{Mod},
&M&\longmapsto P^*\otimes_EM
\end{aligned}
\]
互为拟逆。\(G\) 也自然同构于 \(\operatorname{Hom}_E(P,-)\)，其中该 Hom 模的左 \(R\) 作用为 \((rf)(p)=f(pr)\)。
\end{proposition}
\begin{proof}
对有限投射模，对偶两次的求值映射 \(P\rightarrow P^{**}\) 是同构：这在有限自由模上由坐标直接成立，限制到直和因子后仍成立。因此取对偶给出
\[
\operatorname{End}_R(P^*)\cong E^{\mathrm{op}}.
\]
对偶模有限投射，而且在每个局部环上的秩与 \(P\) 相同，均为正。若 \(rP^*=0\)，在每个极大理想处用正秩自由模得 \(r/1=0\)，再由局部消失判据应用于 \(Rr\) 得 \(r=0\)。故 \(P^*\) 忠实，由\cref{b2:lem:azumaya-projective-base-change}还是生成元。

将\cref{b2:thm:morita-equivalence}用于左 \(R\) 模 \(P^*\)。该定理的自同态反环现在为 \(E\)，给出等价函子 \(\operatorname{Hom}_R(P^*,-)\) 及拟逆 \(P^*\otimes_E-\)。由\cref{b2:prop:finite-projective-hom-tensor}和双对偶同构，前一个函子自然同构于 \(P\otimes_R-\)，其中左 \(E\) 作用恰为求值作用。

等价的全忠实性又给出关于 \(N\) 自然的同构
\[
\operatorname{Hom}_E(P,P\otimes_RN)
\cong\operatorname{Hom}_R(R,N)\cong N.
\]
对任意左 \(E\) 模 \(M\)，已有自然同构 \(M\cong P\otimes_R(P^*\otimes_EM)\)。在上式取 \(N=P^*\otimes_EM\)，得到 \(\operatorname{Hom}_E(P,M)\cong P^*\otimes_EM\)，因此这个 Hom 函子也给出拟逆。
\end{proof}

以 \(R=k[t]\)、\(P=R^n\)、\(n\ge1\) 为例，取 \(a\in k\)、\(r\ge1\)。矩阵环的模范畴仍记录多项式变量 \(t\) 的作用。模 \(R/(t-a)\) 经上述等价成为 \(k^n\)，其中 \(t\) 作标量 \(a\)；模 \(R/(t-a)^r\) 则成为 \(\bigl(R/(t-a)^r\bigr)^n\)，保留 \((t-a)\) 的幂零层数。矩阵大小改变了底向量空间维数，却没有丢掉这些局部信息。

\ProseParagraphBreak
这与\cref{b2:prop:morita-center}相呼应：模范畴恒等函子的自然自变换环是原环的中心，因而 Morita 等价保持中心。对 Azumaya \(R\) 代数，中心由\cref{b2:prop:azumaya-center}确定为 \(R\)，所以模范畴仍保留这个交换基环。另一方面，中心相同不保证 Morita 等价；\(\mathbb R\) 与 \(\mathbb H\) 的中心相同，它们的唯一简单模的自同态除环分别为 \(\mathbb R\) 与 \(\mathbb H^{\mathrm{op}}\)，不可能由范畴等价相互对应。非零 Brauer 类正体现了这种中心之外的代数差别。

\ProseParagraphBreak
本节使用的几何语言都有具体的模论内容：主开集对应局部化，局部基对应向量丛的坐标，矩阵共轭对应自同态代数的粘合，模范畴等价对应 Morita 理论。一般非交换代数的几何还会研究分次模、商范畴及导出范畴等对象。进一步学习这些方向时，需要另行建立相应的几何或同调基础；不能仅凭一个非交换环的中心，就断言它具有某个完整的交换几何模型。
